\section{Introduction}
\label{sec:intro}

Language models encode rich representations of human personas in their activations. Recent work has shown that simple linear directions in a model's residual stream can capture personality traits~\citep{chen2025persona}, ethical dispositions~\citep{zou2023representation}, and professional roles~\citep{poterti2025designing}, and that adding or subtracting these directions at inference time steers model behavior accordingly~\citep{turner2023activation}. This finding---that personas are linearly represented---has opened a practical toolkit for controlling LLM outputs without fine-tuning.

Yet each persona vector is typically extracted in isolation: one direction for agreeableness, another for political conservatism, another for Buddhism. A natural question follows: {\bf do these vectors share underlying structure?} If persona representations are not independent but instead lie in a low-dimensional subspace, then a small set of basis components could span the full diversity of personas. Discovering such a basis would reveal the ``atoms'' of persona as encoded by the model, enable principled composition of arbitrary personas from a compact set of directions, and test whether the model's internal geometry mirrors known psychological structures like the Big Five~\citep{goldberg1990alternative}.

Prior work has begun to address pieces of this question. \citet{feng2026persona} show that Big Five trait vectors are approximately orthogonal and support algebraic composition, but they study only five predefined dimensions. \citet{suh2024rediscovering} recover Big Five factors via SVD of output probabilities, but operate in probability space rather than internal representations. \citet{cintas2025localizing} localize persona representations across layers using PCA, but analyze each persona separately rather than decomposing the joint space. No study has extracted vectors for a large, diverse collection of personas and applied PCA to discover the natural dimensionality and primary components of the full persona space.

We address this gap directly. We extract contrastive persona vectors for \npersonas diverse personas from \qwen across multiple layers of the residual stream, stack them into a matrix $\mV \in \sR^{104 \times 3584}$, and apply PCA to discover the principal axes of variation. Our analysis reveals that persona space has an effective dimensionality of \effdim---far less than the 104 individual vectors---and that the top 20 principal components are all statistically significant by permutation testing. The components are interpretable: PC1 captures a nihilism-versus-agency axis, PC2 separates antisocial traits from cultural interests, and later components pick out political, ethical, and religious dimensions. Critically, these axes cross-cut standard psychological categories; the Big Five does not dominate the decomposition. We validate the causal relevance of the discovered components via activation steering, showing that injecting PC1 shifts model behavior on nihilism-related statements from 10\% to 40\% agreement.

Our contributions are:
\begin{itemize}[leftmargin=*,itemsep=0pt,topsep=0pt]
    \item We perform the first PCA decomposition of a large, diverse collection of persona vectors (\npersonas personas spanning 10 semantic categories) extracted from LLM residual streams, finding an effective dimensionality of \effdim.
    \item We show that the model's natural persona basis does not align with human-defined taxonomies such as the Big Five---the principal components cross-cut semantic categories, revealing a distinct internal geometry.
    \item We validate the causal relevance of discovered principal components through activation steering experiments, confirming that basis directions modulate model behavior.
    \item We characterize how persona structure varies across layers, finding that later-middle layers (layer 20 of 28) concentrate the most structured persona representations.
\end{itemize}
